\documentclass[11pt]{article}
\usepackage[utf8]{inputenc}
\usepackage[T1]{fontenc}
\usepackage[spanish]{babel}
\usepackage{amsmath,amssymb}
\usepackage{geometry}
\usepackage{booktabs}
\usepackage{hyperref}
\usepackage{xcolor}
\usepackage{graphicx}
\usepackage{tikz}
\usetikzlibrary{arrows.meta,positioning,fit,shapes.geometric,shapes.misc,calc}

\geometry{letterpaper, margin=1in}
\hypersetup{colorlinks=true, linkcolor=blue!45!black, urlcolor=blue!45!black}

\definecolor{cInput}{HTML}{E0F2FE}
\definecolor{cRun}{HTML}{DCFCE7}
\definecolor{cData}{HTML}{FEF3C7}
\definecolor{cAnalysis}{HTML}{EDE9FE}
\definecolor{cWarn}{HTML}{FEE2E2}
\definecolor{cGraph}{HTML}{CCFBF1}
\definecolor{cOpt}{HTML}{FCE7F3}

\tikzset{
  block/.style={rectangle, rounded corners=2mm, draw=black!55, very thick, align=center, minimum width=31mm, minimum height=10mm, font=\small},
  smallblock/.style={rectangle, rounded corners=2mm, draw=black!55, thick, align=center, minimum width=25mm, minimum height=8mm, font=\scriptsize},
  phase/.style={rectangle, rounded corners=2mm, draw=black!55, very thick, align=center, minimum width=27mm, minimum height=13mm, font=\scriptsize},
  decision/.style={diamond, aspect=2.2, draw=black!55, very thick, align=center, inner sep=1pt, font=\scriptsize},
  arrow/.style={-{Latex[length=2.5mm]}, thick, draw=black!70},
  dashedarrow/.style={-{Latex[length=2.5mm]}, thick, dashed, draw=black!55},
  group/.style={rectangle, rounded corners=3mm, draw=black!35, dashed, inner sep=4mm}
}

\title{Algoritmo local para recolección secuencial y análisis de narrativas}
\author{}
\date{}

\begin{document}
\maketitle

\begin{abstract}
Este documento describe el algoritmo implementado, desde la aplicación local,
para construir y analizar un corpus narrativo trazable. La propuesta separa
rubros de variantes, años y capas de fuente; guarda bases parciales; fusiona
registros en un archivo estructurado; limpia los textos antes de interpretarlos;
produce indicios cuantitativos, mapas de relaciones narrativas, disección
estructural de proposiciones y selecciones revisables de actores, ideas o
fuentes que merecen lectura cercana. El diseño evita mezclar noticias, foros y
artículos científicos como si fueran una única voz homogénea. La tesis
metodológica es pasar de \textit{Big Data} a \textit{Smart Data}: no importa
acumular textos indiscriminadamente, sino construir evidencia pública,
estratificada, trazable y conceptualmente pertinente. El resultado no es una
interpretación automática, sino una infraestructura de lectura: un corpus
organizado, auditable y preparado para contrastar hipótesis narrativas.
\end{abstract}

\section{Principio metodológico}

El sistema se diseña para tópicos cambiantes. El tópico puede ser
\textit{tatuaje}, \textit{programación con IA}, \textit{violencia}, u otro. El
punto crítico es no buscar una cadena larga con conectores como si fuera una
consulta única. En lugar de ello, se separan rubros conceptuales:
\[
R=\{r_1,r_2,\ldots,r_p\},
\]
años:
\[
Y=\{y_0,y_0+1,\ldots,y_f\},
\]
y capas de fuente:
\[
S=\{\text{noticias},\text{foros},\text{artículos},\text{reportes}\}.
\]

La corrida publicable se ejecuta de manera lineal:
\[
\forall r\in R,\quad \forall y\in Y,\quad \forall s\in S.
\]
Así se sabe qué parte terminó, qué parte falló y qué capa quedó vacía. Esto es
preferible a una búsqueda mezclada que produce una base grande pero sesgada.

Este procedimiento no busca medir directamente una realidad social, sino
construir un corpus trazable de huellas discursivas. Su utilidad está en
permitir comparar cómo distintos géneros de fuente ---prensa, foros, documentos
institucionales y literatura académica--- estabilizan, transforman o disputan
una narrativa a lo largo del tiempo.

La unidad de análisis no es sólo la palabra clave. El sistema identifica
fuentes, marcos argumentales, cambios temporales, silencios relevantes,
proposiciones y relaciones entre actores, ideas, documentos y capas de fuente.
Por ello, los resultados deben leerse como señales estructurales para guiar una
lectura crítica, no como veredictos de verdad ni como prueba judicial cerrada.

\section{Fases del procedimiento}

Para facilitar la lectura desde humanidades, el algoritmo puede entenderse como
nueve fases. Cada fase responde una pregunta metodológica distinta. La parte
computacional no sustituye la interpretación; organiza materiales, deja rastro
de las decisiones y permite comparar capas discursivas.

\begin{table}[ht]
\centering
\small
\begin{tabular}{p{0.16\textwidth}p{0.37\textwidth}p{0.37\textwidth}}
\toprule
\textbf{Fase} & \textbf{Qué se hace} & \textbf{Por qué se hace} \\
\midrule
1. Delimitar & Se define el tópico, la región, los años y los rubros de búsqueda. &
Evita que el corpus sea una mezcla accidental de sentidos distintos del mismo término. \\
\addlinespace
2. Reunir voces & Se buscan documentos por año y por capa: noticias, foros,
artículos y reportes. & Permite distinguir diálogo social, prensa, producción
académica y documentos institucionales. \\
\addlinespace
3. Depurar & Se limpian textos, se eliminan duplicados, menús, publicidad,
palabras vacías y exclusiones conceptuales. & Reduce ruido antes de construir
patrones; no se interpreta texto contaminado como si fuera discurso social. \\
\addlinespace
4. Describir & Se calculan frecuencias, frases repetidas, tonalidad léxica
exploratoria y eventos narrativos. & Produce una primera lectura cuantitativa y
auditable del corpus. \\
\addlinespace
5. Podar inteligentemente & Se identifican fuentes con mayor peso estructural,
marcos repetidos y documentos que sólo actúan como eco. & Evita confundir
volumen documental con diversidad narrativa. \\
\addlinespace
6. Diseccionar & Se extraen proposiciones, actos de habla, relaciones causales,
marcadores retóricos revisables y ausencias relativas dentro del corpus. & Permite
pasar de contar palabras a observar estructura argumental. \\
\addlinespace
7. Relacionar & Se construyen redes de documentos, actores, fuentes, años,
conceptos, proposiciones y etapas narrativas. & Permite observar cómo se
conectan voces, temas y momentos del relato. \\
\addlinespace
8. Seleccionar & Se resuelve el problema del cubridor nodal sobre la red
ponderada y se comparan métodos bajo el mismo presupuesto. & Resume la red en
un conjunto revisable de nodos que cubren relaciones relevantes. \\
\addlinespace
9. Sintetizar & Se exportan corpus, tablas, grafos, trazabilidad técnica y
resultados de métodos. & Facilita comparar soluciones, revisar evidencia y
preparar análisis interpretativo o publicación. \\
\bottomrule
\end{tabular}
\caption{Fases del sistema explicadas para lectores no especializados en computación.}
\label{fig:tabla-fases}
\end{table}

\begin{figure}[ht]
\centering
\resizebox{\textwidth}{!}{%
\begin{tikzpicture}[node distance=7mm and 5mm]
  \node[phase, fill=cInput] (f1) {1\\Delimitar\\tema y periodo};
  \node[phase, fill=cRun, right=of f1] (f2) {2\\Reunir\\voces};
  \node[phase, fill=cData, right=of f2] (f3) {3\\Depurar\\textos};
  \node[phase, fill=cAnalysis, right=of f3] (f4) {4\\Describir\\patrones};
  \node[phase, fill=cData, below=of f4] (f5) {5\\Podar\\Smart Data};
  \node[phase, fill=cAnalysis, left=of f5] (f6) {6\\Diseccionar\\proposiciones};
  \node[phase, fill=cGraph, left=of f6] (f7) {7\\Relacionar\\en redes};
  \node[phase, fill=cOpt, left=of f7] (f8) {8\\Seleccionar\\cubridor};
  \node[phase, fill=cWarn, below=of f8] (f9) {9\\Sintetizar\\evidencia};

  \draw[arrow] (f1) -- (f2);
  \draw[arrow] (f2) -- (f3);
  \draw[arrow] (f3) -- (f4);
  \draw[arrow] (f4) -- (f5);
  \draw[arrow] (f5) -- (f6);
  \draw[arrow] (f6) -- (f7);
  \draw[arrow] (f7) -- (f8);
  \draw[arrow] (f8) -- (f9);
  \draw[dashedarrow] (f9.south) .. controls +(0,-16mm) and +(0,-16mm) .. node[below,font=\scriptsize] {revisión humana y ajuste del corpus} (f1.south);
\end{tikzpicture}%
}
\caption{Resumen de fases para lectura humanística: el proceso es iterativo y auditable.}
\label{fig:fases}
\end{figure}

Las fases no constituyen una cadena mecánica de descubrimiento. Cada una reduce
una forma distinta de ambigüedad: la delimitación reduce ambigüedad conceptual;
la recolección reduce ambigüedad de procedencia; la depuración reduce ruido
documental; la descripción produce indicios; la poda evita acumulación
redundante; la disección separa proposiciones y actos de habla; la red muestra
relaciones posibles; el cubridor resume estructura; y la síntesis prepara
preguntas interpretativas que deben volver al corpus.

\section{Arquitectura general}

\begin{figure}[ht]
\centering
\resizebox{\textwidth}{!}{%
\begin{tikzpicture}[node distance=10mm and 9mm]
  \node[block, fill=cInput] (topic) {Pregunta de estudio\\tópico y sentido};
  \node[block, fill=cInput, right=of topic] (rubrics) {Rubros de búsqueda\\núcleo, identidad,\ salud, oficio};
  \node[block, fill=cInput, right=of rubrics] (scope) {Marco del corpus\\región y años};

  \node[block, fill=cRun, below=of topic] (plan) {Plan de lectura\\rubro $\times$ año $\times$ fuente};
  \node[block, fill=cRun, right=of plan] (search) {Reunir voces\\prensa, foros,\ artículos, reportes};
  \node[block, fill=cRun, right=of search] (fetch) {Recuperar material\\texto completo o señal parcial};

  \node[block, fill=cData, below=of plan] (clean) {Depurar textos\\limpieza y exclusiones};
  \node[block, fill=cData, right=of clean] (partial) {Guardar rastro\\una base por paso};
  \node[block, fill=cData, right=of partial] (merge) {Unir corpus\\sin duplicados};

  \node[block, fill=cAnalysis, below=of clean] (analysis) {Describir relatos\\frases, actores,\ eventos, tono};
  \node[block, fill=cGraph, right=of analysis] (graphs) {Relacionar elementos\\red narrativa y red semántica};
  \node[block, fill=cOpt, right=of graphs] (cover) {Sintetizar estructura\\nodos que explican relaciones};

  \node[block, fill=cData, below=of graphs] (export) {Archivo final\\corpus y análisis reproducible};

  \draw[arrow] (topic) -- (rubrics);
  \draw[arrow] (rubrics) -- (scope);
  \draw[arrow] (topic) -- (plan);
  \draw[arrow] (rubrics) -- (plan);
  \draw[arrow] (scope) -- (plan);
  \draw[arrow] (plan) -- (search);
  \draw[arrow] (search) -- (fetch);
  \draw[arrow] (fetch) -- (clean);
  \draw[arrow] (clean) -- (partial);
  \draw[arrow] (partial) -- (merge);
  \draw[arrow] (merge) -- (analysis);
  \draw[arrow] (analysis) -- (graphs);
  \draw[arrow] (graphs) -- (cover);
  \draw[arrow] (analysis) -- (export);
  \draw[arrow] (graphs) -- (export);
  \draw[arrow] (cover) -- (export);
\end{tikzpicture}%
}
\caption{Arquitectura general expresada como flujo de decisiones interpretativas y cómputo local.}
\label{fig:arquitectura}
\end{figure}

La Figura~\ref{fig:arquitectura} debe leerse de izquierda a derecha y de arriba
abajo. La primera línea corresponde a decisiones de investigación: qué se estudia, con qué
variantes y en qué marco temporal o geográfico. La segunda línea corresponde a
la recolección: reunir voces sin confundirlas. La tercera línea corresponde al
cuidado del corpus: limpiar, registrar y fusionar sin perder trazabilidad. La
última línea corresponde al análisis: describir patrones, observar relaciones y
producir una síntesis revisable.

\section{Corrida secuencial por rubro, término, año y fuente}

La recolección se define como un conjunto de tareas:
\[
T=\{(r,v,y,s): r\in R,\ v\in V_r,\ y\in Y,\ s\in S^\star\},
\]
donde \(V_r\) es la lista ordenada de términos y sinónimos del rubro \(r\), y
\(S^\star\subseteq S\) son las capas seleccionadas por el usuario. Cada tarea
usa un solo término \(v\), no todos los sinónimos a la vez. Si una cuota anual
por tipo de fuente ya se cumplió, los pasos restantes de ese año y capa se
saltan. Cada tarea produce una base parcial \(D_{r,v,y,s}\). La base total antes
de análisis es:
\[
D=\bigcup_{r\in R}\bigcup_{v\in V_r}\bigcup_{y\in Y}\bigcup_{s\in S^\star}D_{r,v,y,s},
\]
desduplicada por URL o título.

\begin{figure}[ht]
\centering
\begin{tikzpicture}[node distance=9mm and 7mm]
  \node[smallblock, fill=cInput] (start) {Inicio};
  \node[smallblock, fill=cRun, below=of start] (rubro) {Elegir un rubro\\del tema};
  \node[smallblock, fill=cRun, below=of rubro] (term) {Elegir un término\\del rubro};
  \node[smallblock, fill=cRun, below=of term] (year) {Elegir un año};
  \node[smallblock, fill=cRun, below=of year] (source) {Elegir una capa\\de fuente};
  \node[smallblock, fill=cRun, below=of source] (query) {Formular búsqueda\\clara y breve};
  \node[smallblock, fill=cRun, below=of query] (search) {Consultar\\fuentes públicas};
  \node[decision, fill=cWarn, below=of search] (found) {¿hay\\resultados?};
  \node[smallblock, fill=cData, below left=10mm and 13mm of found] (empty) {Registrar\\vacío o falla};
  \node[smallblock, fill=cRun, below right=10mm and 13mm of found] (download) {Recuperar texto\\o resumen};
  \node[smallblock, fill=cData, below=of download] (save) {Guardar\\evidencia};
  \node[decision, fill=cWarn, below=of save] (more) {¿faltan\\pasos?};
  \node[smallblock, fill=cData, below=of more] (merge) {Unir corpus\\sin duplicados};

  \draw[arrow] (start) -- (rubro);
  \draw[arrow] (rubro) -- (term);
  \draw[arrow] (term) -- (year);
  \draw[arrow] (year) -- (source);
  \draw[arrow] (source) -- (query);
  \draw[arrow] (query) -- (search);
  \draw[arrow] (search) -- (found);
  \draw[arrow] (found) -- node[left,font=\scriptsize] {no} (empty);
  \draw[arrow] (found) -- node[right,font=\scriptsize] {sí} (download);
  \draw[arrow] (download) -- (save);
  \draw[arrow] (save) -- (more);
  \draw[arrow] (empty) |- (more);
  \draw[arrow] (more) -- node[right,font=\scriptsize] {no} (merge);
  \draw[dashedarrow] (more.west) .. controls +(-34mm,0) and +(-34mm,0) .. node[left,font=\scriptsize] {sí} (term.west);
\end{tikzpicture}
\caption{Ejecución secuencial adaptativa. Cada paso reporta rubro, término, año, capa de fuente y número de documentos.}
\end{figure}

Esta fase existe para evitar una ilusión frecuente: obtener muchos documentos
pero de una sola clase. Si sólo aparecen artículos científicos, no hay base
suficiente para hablar de conversación social; si sólo aparecen noticias, falta
contrastar con foros o producción académica. La corrida secuencial muestra los
vacíos en lugar de ocultarlos.

\section{Extracción responsable y estados de evidencia}

La recolección distingue entre encontrar una referencia, descargar texto completo
y conservar sólo una señal parcial. Esta distinción es ética y metodológica:
evita tratar un resumen RSS, un título indexado o una página con restricción
como si fuera texto completo.

\begin{figure}[ht]
\centering
\resizebox{0.98\textwidth}{!}{%
\begin{tikzpicture}[node distance=8mm and 8mm]
  \node[block, fill=cInput] (index) {Índice público\\GDELT, RSS,\ OpenAlex, Crossref};
  \node[block, fill=cRun, right=of index] (candidate) {URL o metadato\\con fuente, año,\ rubro y término};
  \node[decision, fill=cWarn, right=of candidate] (allowed) {¿permite\\robots.txt?};
  \node[block, fill=cRun, below right=8mm and 8mm of allowed] (full) {Descargar HTML/PDF\\si es público};
  \node[block, fill=cData, above right=8mm and 8mm of allowed] (partial) {Guardar señal parcial\\título/resumen/metadato};
  \node[block, fill=cData, right=of full] (status) {Estado explícito\\ok, ok\_partial,\ too\_short, error};
  \node[block, fill=cAnalysis, right=of partial] (audit) {Auditoría\\acceso, fuente,\ permiso, fecha};
  \node[block, fill=cGraph, right=of status] (analysis) {Análisis local\\sin subir textos\\a servicios externos};

  \draw[arrow] (index) -- (candidate);
  \draw[arrow] (candidate) -- (allowed);
  \draw[arrow] (allowed) -- node[above,font=\scriptsize] {no/paywall/parcial} (partial);
  \draw[arrow] (allowed) -- node[right,font=\scriptsize] {sí} (full);
  \draw[arrow] (full) -- (status);
  \draw[arrow] (partial) -- (audit);
  \draw[arrow] (audit) |- (status);
  \draw[arrow] (status) -- (analysis);
\end{tikzpicture}%
}
\caption{Extracción responsable: el sistema no fuerza sitios cerrados; conserva metadatos parciales cuando no hay permiso de texto completo.}
\label{fig:extraccion-responsable}
\end{figure}

En consecuencia, los estados no son detalles técnicos sino categorías de
evidencia. \texttt{ok} indica texto suficiente para análisis textual completo;
\texttt{ok\_partial} indica una señal pública útil pero limitada; \texttt{too\_short}
preserva registros breves que no deben inflar análisis de texto; y los errores
se reportan como límites de la muestra. La publicación debe informar cuántos
registros pertenecen a cada estado.

\section{Limpieza, composición de frases y análisis local}

Para cada documento \(d\), el texto se normaliza antes de construir conceptos:
\[
\text{text}(d)\to \text{lowercase}\to \text{normalización}\to
\text{stopwords}\to \text{tokens}.
\]
Luego se extraen monogramas, bigramas y trigramas. Si una frase repetida explica
la mayor parte de las apariciones de sus componentes, se usa como nodo canónico:
\[
q\prec p
\quad\text{si}\quad
q\subset p,\qquad
\frac{freq(p)}{freq(q)}\geq \theta,\qquad freq(p)\geq 2.
\]

\begin{figure}[ht]
\centering
\begin{tikzpicture}[node distance=10mm and 8mm]
  \node[block, fill=cData] (raw) {Material reunido\\notas, foros,\ artículos};
  \node[block, fill=cData, right=of raw] (norm) {Texto ordenado\\sin menús,\ ruido ni duplicados};
  \node[block, fill=cData, right=of norm] (tokens) {Palabras útiles\\vocabulario depurado};

  \node[block, fill=cAnalysis, below=of raw] (ngrams) {Expresiones frecuentes\\una, dos o tres palabras};
  \node[block, fill=cAnalysis, right=of ngrams] (phrases) {Frases estables\\conceptos repetidos};
  \node[block, fill=cAnalysis, right=of phrases] (events) {Relato por documento\\actores, conflicto,\ resolución};

  \node[block, fill=cGraph, below=of phrases] (graphs) {Mapas de relaciones\\frecuencia y conexión};

  \draw[arrow] (raw) -- (norm);
  \draw[arrow] (norm) -- (tokens);
  \draw[arrow] (tokens) -- (ngrams);
  \draw[arrow] (ngrams) -- (phrases);
  \draw[arrow] (tokens) -- (events);
  \draw[arrow] (phrases) -- (graphs);
  \draw[arrow] (events) -- (graphs);
\end{tikzpicture}
\caption{Preparación del texto antes de construir mapas de relaciones.}
\end{figure}

La limpieza no es un trámite técnico menor. En páginas web, el texto puede traer
menús, publicidad, biografías de autores, notas relacionadas o palabras vacías.
Si ese ruido entra al análisis, la red puede destacar términos irrelevantes. Por
eso se depura antes de contar palabras, detectar frases o construir relaciones.
La composición de frases busca respetar unidades de sentido: una expresión
repetida como \textit{arte corporal} o \textit{tatuaje religioso} puede ser más
informativa que sus palabras separadas.

\section{Tonalidad léxica local y radar por fuente}

La tonalidad léxica se calcula con un léxico local bilingüe. Para cada documento
se parte de:
\[
tono(d)=
\frac{pos(d)-neg(d)}{pos(d)+neg(d)},
\]
con \(tono(d)=0\) si no hay términos positivos ni negativos. En la aplicación
se incorporan reglas simples de negación e intensificación: expresiones como
``no seguro'' invierten la orientación inmediata, mientras que intensificadores
como ``muy'' aumentan el peso local de una palabra valorativa. Para graficar en
radar se usa:
\[
rad(d)=\frac{tono(d)+1}{2}\in[0,1].
\]
El promedio por año y fuente es:
\[
\bar{s}_{y,c}=
\frac{1}{|D_{y,c}|}\sum_{d\in D_{y,c}}tono(d).
\]
Cada eje del radar representa un tipo de fuente. Esta medición es exploratoria:
sirve para comparar capas discursivas, no para inferir emociones individuales.

En términos interpretativos, el radar no dice ``lo que siente la sociedad'' ni
describe estados psicológicos de autores individuales. Dice si, dentro del
corpus recolectado, una capa discursiva usa con mayor frecuencia vocabulario
valorativo positivo, negativo o mixto. Por eso debe leerse junto con ejemplos
textuales y con la distribución de fuentes por año.

\section{Redes construidas}

\begin{figure}[ht]
\centering
\begin{tikzpicture}[node distance=18mm]
  \node[circle, draw, fill=cData, minimum size=15mm] (doc) {Texto};
  \node[circle, draw, fill=cAnalysis, above left=of doc] (actor) {Actor};
  \node[circle, draw, fill=cAnalysis, above right=of doc, align=center] (stage) {Momento\\relato};
  \node[circle, draw, fill=cGraph, right=of doc] (concept) {Idea};
  \node[circle, draw, fill=cInput, below left=of doc] (source) {Fuente};
  \node[circle, draw, fill=cInput, below=of doc] (year) {Año};
  \node[circle, draw, fill=cInput, below right=of doc] (loc) {Lugar};

  \draw[arrow] (doc) -- node[font=\scriptsize,left] {menciona} (actor);
  \draw[arrow] (doc) -- node[font=\scriptsize,right] {contiene} (stage);
  \draw[arrow] (doc) -- node[font=\scriptsize,above] {usa} (concept);
  \draw[arrow] (source) -- node[font=\scriptsize,left] {publica} (doc);
  \draw[arrow] (year) -- node[font=\scriptsize] {fecha} (doc);
  \draw[arrow] (loc) -- node[font=\scriptsize,right] {localiza} (doc);
  \draw[arrow] (actor) -- node[font=\scriptsize,above] {participa en} (stage);
  \draw[arrow] (stage) -- node[font=\scriptsize,above] {asocia} (concept);
\end{tikzpicture}
\caption{Red narrativa: cada texto conecta actores, ideas, momentos del relato, fuentes, años y lugares.}
\label{fig:red-narrativa}
\end{figure}

La red se usa como mapa de relaciones, no como sustituto de lectura crítica.
Permite observar qué actores aparecen cerca de qué ideas, qué momentos del
relato se repiten, qué fuentes concentran ciertos temas y cómo cambian esos
patrones por año o localización.

Los pesos se asignan por conteo:
\[
w_a=\sum_{d\in D}\mathbf{1}\{a \text{ aparece en }d\}.
\]
Se conserva el peso crudo y se normaliza:
\[
\widetilde{w}_a=\frac{w_a}{\max_{b\in A}w_b}.
\]
La red semántica usa coocurrencias entre conceptos. El grafo de conocimiento
agrega conceptos, fuentes, años, idioma, localización y tipo de fuente. Cuando
Louvain está disponible, se calculan comunidades; si no, se reporta la
alternativa local usada.

\section{Cubridor nodal multiobjetivo}

Esta sección formaliza una pregunta interpretativa simple: si el mapa narrativo
contiene demasiados elementos para revisarlos uno por uno, ¿qué conjunto pequeño
de actores, ideas, fuentes o momentos del relato permite iniciar una lectura
crítica sin borrar la complejidad de la red? En adelante llamamos nodos
relevantes a los elementos que concentran relaciones dentro del corpus. La
formalización matemática permite comparar distintas selecciones, pero la
relevancia final depende de lectura crítica y contraste con ejemplos textuales.

Sea \(G=(V,A)\) una red ponderada. El universo del problema es el conjunto de
aristas objetivo:
\[
{\mathcal U}=A^\star\subseteq A.
\]
Cada nodo candidato \(v\in B\subseteq V\) cubre sus aristas incidentes:
\[
S_v=\{a\in A^\star:a\sim v\}.
\]
La solución es un conjunto \(C\subseteq B\). Con variables binarias \(x_v\) y
\(r_a\):
\[
x_v=1\Leftrightarrow v\in C,\qquad
r_a=1\Leftrightarrow a\text{ toca algún nodo seleccionado}.
\]
La condición de arista removida se modela con:
\[
r_a\geq b_{av}x_v,\quad \forall a\in A^\star,\ \forall v\in B,
\]
\[
r_a\leq\sum_{v\in B}b_{av}x_v,\quad \forall a\in A^\star.
\]
Esto evita confundir el daño estructural con el subgrafo inducido. Una arista se
cuenta como perdida si al retirar el conjunto seleccionado toca al menos un nodo
del cubridor.

Los objetivos normalizados son:
\[
f_1(C)=\frac{\sum_v x_v}{|B|}\quad \text{minimizar},
\]
\[
f_2(C)=\frac{\sum_v\sigma_vx_v}{\sum_v\sigma_v}\quad \text{maximizar},
\]
\[
f_3(C)=\frac{\sum_a w_ar_a}{\sum_a w_a}\quad \text{minimizar}.
\]

\begin{figure}[ht]
\centering
\resizebox{0.96\textwidth}{!}{%
\begin{tikzpicture}[node distance=10mm and 9mm]
  \node[block, fill=cData] (corpus) {Corpus depurado\\textos y procedencia};
  \node[block, fill=cGraph, right=of corpus] (graph) {Mapa de relaciones\\actores, ideas,\ fuentes};
  \node[block, fill=cOpt, right=of graph] (candidates) {Nodos candidatos\\posibles síntesis};
  \node[block, fill=cOpt, right=of candidates] (archive) {Opciones no superadas\\varias síntesis posibles};
  \node[block, fill=cAnalysis, below=of archive] (review) {Revisión humana\\volver a ejemplos\ textuales};
  \node[block, fill=cAnalysis, below=of graph] (metrics) {Comparación común\\tamaño, relevancia,\ daño estructural};

  \draw[arrow] (corpus) -- (graph);
  \draw[arrow] (graph) -- (candidates);
  \draw[arrow] (candidates) -- (archive);
  \draw[arrow] (archive) -- (review);
  \draw[arrow] (archive) -- (metrics);
  \draw[dashedarrow] (review) .. controls +(0,-16mm) and +(0,-16mm) .. node[below,font=\scriptsize] {ajuste interpretativo} (corpus);
\end{tikzpicture}%
}
\caption{Síntesis revisable: del corpus al mapa, del mapa a nodos relevantes y de allí a lectura crítica.}
\label{fig:sintesis}
\end{figure}

La última fase selecciona un conjunto pequeño de nodos que ayude a explicar el
mapa sin destruirlo. Esta selección no pretende reemplazar el análisis
interpretativo. Funciona como una herramienta para preguntar: ¿qué actores,
ideas o fuentes son centrales?, ¿qué relaciones se perderían si quitáramos esos
nodos?, ¿hay varias síntesis igualmente defendibles?

La aplicación puede comparar varias estrategias de búsqueda de soluciones:
barrido glotón por prioridades, búsqueda evolutiva multiobjetivo, recocido
multiobjetivo y una adaptación multiobjetivo del método de composición musical.
Para no cargar la figura con siglas, la Figura~\ref{fig:sintesis} muestra la
pregunta común que comparten esos métodos.

La factibilidad se compara con criterio de Coello:
\[
\text{factible} \succ \text{infactible}.
\]
Entre dos infactibles gana la menor violación total. Entre dos factibles se
aplica dominancia de Pareto. Todos los métodos deben usar el mismo presupuesto
de evaluaciones de la función objetivo.

El MMC-MO resuelve una sola vez las relajaciones monoobjetivo de referencia y
las guarda como memoria inicial. Después modifica esa guía con soluciones no
dominadas encontradas durante la búsqueda; no debe resolver los problemas
relajados en cada iteración.

\section{Exportación final}

La salida de análisis se integra en:
\[
\texttt{narrative\_analysis\_unified.json}.
\]
Este archivo contiene registros filtrados, tonalidad léxica, eventos narrativos,
actores, grupos de ideas, red narrativa, red semántica, grafo de conocimiento,
cubridor y frente Pareto. Es la unidad reproducible para publicación y análisis
posterior.

\end{document}
